\subsection{Fragestellungen}
\label{subsec:fragestellungen}

Sechs Fragen strukturieren Untersuchung und Bewertung:

\begin{enumerate}
  \item Welche fachlich-technischen, qualitativen und betrieblichen
    Anforderungen werden durch das Technologie-Template adressiert?
  \item Wie unterstützen die Architektur und das Profilmodell die
    Bereitstellung unterschiedlicher Varianten für Web-, Cloud- und
    Desktop-Projekte?
  \item Inwiefern fördern der gemeinsame technische Kern und die
    Single-Repository-Strategie die Wiederverwendung und Standardisierung
    zwischen Projekten?
  \item Welches Potenzial bietet das Template, wiederkehrende Aufwände bei
    Projektinitialisierung, Entwicklungsworkflow, automatisierter Codequalitäts-
    und Architekturprüfung sowie Bereitstellung zu reduzieren?
  \item Für welche Projektarten und Kundenanforderungen stellt das Template
    eine geeignete technische Ausgangsbasis dar?
  \item Welche technischen Grenzen, Wartungsaufwände, Risiken und externen
    Abhängigkeiten sind beim Einsatz einschließlich der Governance zu
    berücksichtigen?
\end{enumerate}

Ihre Beantwortung stützt sich auf Anforderungen, Architektur und vorhandene
Verifikationsnachweise. Als querschnittlicher Evaluationsaspekt wird dabei
geprüft, wie das profilbasierte Full-Stack-Template ausgewählte Codequalitäts-
und Architekturregeln sprachübergreifend, konfigurierbar und CI-wirksam
durchsetzt, ohne sie in die Produktlaufzeit zu verlagern. Dieser Aspekt ergänzt
die sechs Fragen und wird nicht als eigenständige, empirisch generalisierbare
Hauptforschungsfrage behandelt.
